\section{第一次作业}

\subsection{1. 应用冲激信号的抽样特性求函数值}
\noindent\textbf{原题：}应用冲激信号的抽样特性（筛选特性），求下列表示式的函数值：
\begin{enumerate}[label=(\arabic*)]
    \item $\displaystyle \int_{-\infty}^{\infty} f(t-t_0)\delta(t)dt$
    \item $\displaystyle \int_{-\infty}^{\infty} f(t_0-t)\delta(t)dt$
    \item $\displaystyle \int_{-\infty}^{\infty} \delta(t-t_0)u(t-2t_0)dt$
    \item $\displaystyle \int_{-\infty}^{\infty} (t+\sin t)\delta\left(t-\frac{\pi}{6}\right)dt$
    \item $\displaystyle \int_{-\infty}^{\infty} e^{-j\omega t}\left[\delta(t)-\delta(t-t_0)\right]dt$
\end{enumerate}

\subsubsection{解答}
冲激信号抽样（筛选）性质：
\[
\int_{-\infty}^{\infty} f(\tau)\delta(\tau-t_0)d\tau = f(t_0)
\]

\begin{enumerate}[label=(\arabic*)]
    \item
    \[
    \int_{-\infty}^{\infty} f(t-t_0)\delta(t)dt
    = f(0-t_0) = \boldsymbol{f(-t_0)}
    \]

    \item
    \[
    \int_{-\infty}^{\infty} f(t_0-t)\delta(t)dt
    = f(t_0-0) = \boldsymbol{f(t_0)}
    \]

    \item
    \[
    \int_{-\infty}^{\infty} \delta(t-t_0)u(t-2t_0)dt
    = u(t_0-2t_0) = u(-t_0)
    \]
    \[
    = \begin{cases}
    \boldsymbol{0}, & t_0>0 \\
    \boldsymbol{1}, & t_0<0 \\
    \text{无定义}, & t_0=0
    \end{cases}
    \]

    \item
    \[
    \int_{-\infty}^{\infty} (t+\sin t)\delta\left(t-\frac{\pi}{6}\right)dt
    = \frac{\pi}{6} + \sin\frac{\pi}{6}
    = \boldsymbol{\frac{\pi}{6}+\frac{1}{2}}
    \]

    \item
    \[
    \int_{-\infty}^{\infty} e^{-j\omega t}\left[\delta(t)-\delta(t-t_0)\right]dt
    = e^{0} - e^{-j\omega t_0}
    = \boldsymbol{1 - e^{-j\omega t_0}}
    \]
\end{enumerate}

\vspace{1em}
\hrule
\vspace{1em}

\subsection{2. 判断周期信号并计算基波周期}
\noindent\textbf{原题：}判断信号 $f(t) = 2\cos(10t+5) - \sin(6t-3)$ 是否为周期信号？如是周期信号，计算基波周期。

\subsubsection{解答}
设两个分量角频率分别为：
\[
\omega_1=10,\quad \omega_2=6
\]
各自周期：
\[
T_1 = \frac{2\pi}{\omega_1} = \frac{\pi}{5},\quad
T_2 = \frac{2\pi}{\omega_2} = \frac{\pi}{3}
\]
周期之比：
\[
\frac{T_1}{T_2} = \frac{3}{5} \in \mathbb{Q}
\]
为有理数，故 $f(t)$ 为**周期信号**。

基波周期为 $T_1,T_2$ 的最小公倍数：
\[
T = \text{LCM}\left(\frac{\pi}{5},\frac{\pi}{3}\right) = \pi
\]

结论：$f(t)$ 是周期信号，基波周期
\[
\boldsymbol{T=\pi}
\]

\vspace{1em}
\hrule
\vspace{1em}

\subsection{3. 绘制卷积波形}
\noindent\textbf{原题：}
$f_1(t)=u(t+1)-u(t-1)$，$f_2(t)=\delta(t+5)+\delta(t-5)$
\begin{enumerate}[label=(\arabic*)]
    \item $s_1(t)=f_1(t)*f_2(t)$
    \item $s_2(t)=\left\{\left[f_1(t)*f_2(t)\right]\left[u(t+5)-u(t-5)\right]\right\}*f_2(t)$
\end{enumerate}

\subsubsection{解答}
利用卷积时移性质：
\[
f(t)*\delta(t-t_0)=f(t-t_0)
\]

\begin{enumerate}[label=(\arabic*)]
    \item
    \[
    s_1(t)
    = f_1(t)*\delta(t+5) + f_1(t)*\delta(t-5)
    = f_1(t+5) + f_1(t-5)
    \]
    其中：
    \[
    f_1(t+5)=u(t+6)-u(t+4),\quad
    f_1(t-5)=u(t-4)-u(t-6)
    \]

    \item
    \[
    \left[f_1(t)*f_2(t)\right]\cdot\big[u(t+5)-u(t-5)\big] = f_1(t+5)
    \]
    \[
    s_2(t)
    = f_1(t+5)*f_2(t)
    = f_1(t+10) + f_1(t)
    \]
\end{enumerate}

\vspace{1em}
\hrule
\vspace{1em}

\subsection{4. 证明正交函数集}
\noindent\textbf{原题：}证明 $\sin t,\sin2t,\dots,\sin nt$ 在 $(0,2\pi)$ 上为正交函数集。
\begin{enumerate}
    \item 是否完备？
    \item 在 $(0,\pi/2)$ 是否正交？
\end{enumerate}

\subsubsection{解答}
正交定义：对任意正整数 $k\neq l$，有
\[
\int_{0}^{2\pi}\sin(kt)\sin(lt)dt=0
\]
利用三角恒等式：
\[
\sin(kt)\sin(lt)=\frac{1}{2}\left[\cos((k-l)t)-\cos((k+l)t)\right]
\]
在 $(0,2\pi)$ 上积分：
\[
\int_{0}^{2\pi}\cos((k-l)t)dt=0,\quad
\int_{0}^{2\pi}\cos((k+l)t)dt=0
\]
故
\[
\int_{0}^{2\pi}\sin(kt)\sin(lt)dt=0,\quad k\neq l
\]
且
\[
\int_{0}^{2\pi}\sin^2(kt)dt=\pi\neq0
\]
因此 $\{\sin nt\}$ 在 $(0,2\pi)$ 上为**正交函数集**。

\begin{enumerate}[label=(\arabic*)]
    \item 该函数集**不完备**。
    只能表示区间上的奇函数，无法表示偶函数，不满足完备性。

    \item 在 $(0,\pi/2)$ 上**不正交**。
    反例：
    \[
    \int_{0}^{\pi/2}\sin t\sin2t\,dt
    = \int_{0}^{\pi/2} 2\sin^2 t\cos t\,dt
    = \frac{2}{3}\neq0
    \]
\end{enumerate}

\vspace{1em}
\hrule
\vspace{1em}

\subsection{5. 计算卷积并绘制波形}
\noindent\textbf{原题：}
已知
\[
f_1(t)=
\begin{cases}
2,&1<t<3\\
0,&\text{其他}
\end{cases},\quad
f_2(t)=
\begin{cases}
1,&0<t<1\\
-1,&1<t<2\\
0,&\text{其他}
\end{cases}
\]
\begin{enumerate}[label=(\arabic*)]
    \item 求 $y(t)=f_1(t)*f_2(t)$，画出概略波形。
    \item 画出 $y(0.5t-2)$ 波形，并说明变换步骤。
\end{enumerate}

\subsubsection{(1) 计算卷积 $y(t) = f_1(t)*f_2(t)$}
卷积定义：
\[
y(t) = \int_{-\infty}^{+\infty} f_1(\tau)\,f_2(t-\tau)\,d\tau
\]
$f_1(\tau)$ 非零区间：$1<\tau<3$
$f_2(t-\tau)$ 非零区间：$t-2<\tau<t$，且
\[
f_2(t-\tau)=
\begin{cases}
1, & t-1<\tau<t\\
-1, & t-2<\tau<t-1
\end{cases}
\]

分区间讨论：
\begin{enumerate}[itemsep=0.5ex]
    \item $t \le 1$：无重叠，$y(t)=0$
    \item $1 < t < 2$：重叠 $1<\tau<t$
    \[
    y(t) = \int_{1}^{t} 2\cdot 1\,d\tau = 2(t-1)
    \]
    \item $2 < t < 3$：重叠分两段
    \[
    y(t) = \int_{1}^{t-1}\!-2\,d\tau + \int_{t-1}^{t}\!2\,d\tau
          = 6-2t
    \]
    \item $3 < t < 4$：重叠 $t-1<\tau<3$
    \[
    y(t) = \int_{t-1}^{3}\!-2\,d\tau = 2t-8
    \]
    \item $t \ge 4$：无重叠，$y(t)=0$
\end{enumerate}

综上：
\[
y(t) =
\begin{cases}
0, & t \leq 1 \\
2(t-1), & 1 < t < 2 \\
6 - 2t, & 2 < t < 3 \\
2t - 8, & 3 < t < 4 \\
0, & t \geq 4
\end{cases}
\]

\subsubsection{(2) 波形变换 $y(0.5t-2)$}
将 $y(0.5t-2)$ 改写为标准形式：
\[
y(0.5t-2) = y\left(0.5(t-4)\right)
\]
变换步骤：
\begin{enumerate}
    \item 原信号 $y(t)$
    \item 右移 4：$y(t-4)$
    \item 横坐标拉伸 2 倍：$y(0.5(t-4))$
\end{enumerate}

原关键点：
\[
t=1,\ 2,\ 3,\ 4
\]
变换后：
\[
t=10,\ 12,\ 14,\ 16
\]

表达式：
\[
y(0.5t-2) =
\begin{cases}
0, & t \leq 10 \\
t-10, & 10 < t < 12 \\
14-t, & 12 < t < 14 \\
t-16, & 14 < t < 16 \\
0, & t \geq 16
\end{cases}
\]